\documentclass{article}

\usepackage{geometry}
\usepackage{comment}
\geometry{
  top=0.75in,
  bottom=0.75in,
  left=0.8in,
  right=0.8in,
  marginparwidth=40pt,
  marginparsep=5pt
}
\usepackage{graphicx}
\usepackage{amsmath}
\usepackage{hyperref}

\title{\textbf{Update \#1: Mobile Starlink Performance over Recurring Route}}

\author{
    \textbf{Jasper McLennan (V01014572)} \\
    \small Department of Computer Science \\
    \small University of Victoria \\
    \small \texttt{jaspermclennan@uvic.ca}
    \and
    \textbf{Rowan Hall (V01041357)} \\
    \small Department of Computer Science \\
    \small University of Victoria \\
    \small \texttt{rowanhall@uvic.ca}
}

\date{\small November 2, 2026}

\begin{document}

\maketitle

%Rowan
\section{Project Overview}

We are studying how a Starlink dish performs on a moving vehicle that drives the same route every day. The dish is mounted on a school bus that runs along Highway~14 between Port Renfrew and the western edge of Greater Victoria. The bus makes the trip twice per school day, once in the morning (around 6:20am) and once in the afternoon (around 2:20pm), and each trip takes about two and a half hours.

A logger on the bus saves data in 30-minute chunks, grouped by trip. The main things it records are:
\begin{itemize}
    \item \textbf{Dish status} (twice a second): latency, packet loss, throughput, signal quality, and whether the dish is obstructed.
    \item \textbf{GPS} (twice a second): the bus's position and speed.
    \item \textbf{Outages}: every time the dish loses connection, with a cause (e.g.\ obstructed, searching for satellites) and a duration.
    \item \textbf{High-rate pings}: a ping every 10~ms, which gives a much finer view of latency than the dish status.
\end{itemize}

We process each chunk with a tool called \texttt{leoviz}, which lines up the GPS with the dish's orientation and estimates which satellite the dish was connected to at each moment. That satellite estimate is what we will use to find handovers.

\section{Progress}

\textbf{Data so far.} We have downloaded and processed 84 trips from September 9 to October 22, 2026 (about 9~GB). Most of these are full morning or afternoon runs, plus some short mid-day recordings.

\textbf{Cleaning pipeline.} We wrote three Python scripts that run in order:
\begin{enumerate}
    \item \textbf{Merge:} combines the dish status, GPS, and outage data for every trip into one table. Each status reading is matched to the closest GPS point in time (within 2 seconds), readings without a GPS match are dropped, and each reading is marked if it happened during an outage.
    \item \textbf{Grid:} splits the route into roughly 100~m squares and calculates the typical latency, throughput, and outage rate in each square.
    \item \textbf{Map:} draws an interactive map of the route coloured by latency, outage rate, or throughput, to put on our project website.
\end{enumerate}

\textbf{First results.} The merged dataset has about 760,000 readings from 57 trips. Some early findings:
\begin{itemize}
    \item Latency is usually good (around 21~ms), but it occasionally spikes to several seconds.
    \item The dish is in an outage about 8\% of the time. The most common cause is the dish searching for satellites, followed by obstructions.
    \item Most of the route has been driven at least 10 times, so we can already see which spots are bad on most trips. About 80 spots (each about 100~m) are in an outage at least 30\% of the time.
    \item Morning and afternoon runs look about the same so far.
\end{itemize}

\textbf{Problems we found.} While checking the data we found some issues that still need fixing:
\begin{itemize}
    \item \textbf{Missing data:} the last chunk of every trip ends with a corrupted line, probably because the logger is cut off when the bus turns off. Our script throws away the whole chunk when this happens, so about a quarter of the chunks (and most short trips) are missing from the dataset.
    \item \textbf{Stopped bus:} about a quarter of the readings were taken while the bus was stopped, which skews the results at bus stops and the depot.
    \item \textbf{Time zones:} times are stored in UTC, so some of our morning/afternoon labels are wrong.
    \item \textbf{Throughput:} no speed tests were run, so throughput only shows how much passengers were using the connection, not how fast it could be.
    \item \textbf{Satellite estimates are noisy:} the estimated satellite often jumps back and forth between candidates, so it needs to be smoothed before we can count handovers.
\end{itemize}


%Jasper
\section{Version Control}

All of our code and reports are in a GitHub repository. The scripts are in \texttt{analysis/}, the project website is in \texttt{gh-pages/}, and these reports are in \texttt{docs/}.

The raw data is too big for Git (over 9~GB), so it is not in the repository. Instead, \texttt{data/readme.txt} has the commands to download it from the shared cloud bucket and process it, and our \texttt{.gitignore} keeps the raw data, zip files, and generated data files out of Git. We learned this the hard way: we accidentally committed a 2.6~GB zip file, GitHub rejected the push, and we had to rewrite our commits to remove it. Small outputs like the final map can still be committed so the website works.

Two small things are still to do: move the processing script out of the ignored data folder so it is tracked, and pin exact package versions in \texttt{requirements.txt}.

\section{Next Steps}
\label{sec:next}

\textbf{Finish cleaning:}
\begin{itemize}
    \item Fix the corrupted-line bug so no chunks are thrown away, then rebuild the dataset.
    \item Download and process the newer trips (up to November 10).
    \item Convert times to local time and label each trip as morning or afternoon and by direction.
    \item Remove or separately label readings taken while the bus was stopped.
    \item Measure position as distance along the route instead of a grid, so trips can be compared side by side.
\end{itemize}

\textbf{Hotspot maps (Weeks 3--4):}
\begin{itemize}
    \item Publish the cleaned map on our website.
    \item For each part of the route, measure how often it has problems across trips, to separate spots that are bad every day (likely trees or hills) from ones that are only bad sometimes.
    \item Compare morning vs.\ afternoon, direction of travel, and week-to-week changes.
\end{itemize}

\textbf{Handover analysis (Weeks 5--6):}
\begin{itemize}
    \item Smooth the satellite estimates and find when the dish switches satellites.
    \item Use the 10~ms ping data to see what happens to latency right around each switch.
    \item Check whether outages happen more often right after a switch, and whether that depends on location.
\end{itemize}

\section{Needed Research}

\begin{itemize}
    \item How \texttt{leoviz} estimates the connected satellite, so we know how much to trust it.
    \item How Starlink schedules satellite switches (other studies report it happens on a fixed 15-second cycle) and how earlier papers measured handovers.
    \item What the dish's outage causes and data fields mean, since Starlink does not officially document them.
    \item Other studies of Starlink on moving vehicles, to compare our results against.
    \item How to match noisy GPS points to a single route line.
    \item Where to get historical weather data and tree/terrain data for the route, to help explain why hotspots happen.
    \item Which methods to use for the final model (e.g.\ anomaly detection, or models that separate location effects from time effects).
\end{itemize}

\section{Problem}
\label{sec:problem}
Low Earth Orbit (LEO) satellite networks like Starlink offer broadband connectivity almost anywhere, but their real world performance changes quite a bit once the terminal starts moving. Mobile terminals are now found mounted on buses, service vehicles, boats, and RVs, where users expect stable internet for navigation, telemetry, or remote work. However, vehicular deployments frequently suffer from sudden latency spikes, poor throughput, and brief connection dropouts.

The underlying causes of these mobile performance drops are not yet fully understood or explored. When a vehicle is in motion, network reliability depends on a mix of overhead obstacles (such as tree canopies, terrain cuts, and overpasses) and rapid orbital dynamics, where terminals must switch beams or satellites every few seconds. Because existing studies look at isolated drives, it is difficult to determine whether a bad connection is an unavoidable spot on the road, an inefficient satellite handover, or bad weather. Understanding where and why these dropouts occur is essential for setting realistic Quality of Service (QoS) expectations, planning routes with reliable connection, and designing smarter handover protocols for mobile satellite systems.


\section{Approach}
\label{sec:approach}

Most mobile Starlink studies rely on a single trip or a handful of drives over different roads. In contrast, our approach uses a dataset collected from a school bus running the exact same route, about two and a half hours each way, twice a day on every school day. Collection began in September 2026 and is still ongoing. The bucket currently holds over 120 sessions (September 9 to November 10), and more trips are added every week.

Because the bus drives the same road over a hundred times across different days, weather, traffic conditions, and use conditions, it acts as a natural controlled experiment. Holding the route constant makes it much easier to figure out what actually causes bad performance, and can indicate whether a drop in service is tied to a specific spot on the road (like trees, hills, or bridges) or to satellite issues overhead.

Our approach focuses on three main steps:

\begin{itemize}
    \item \textbf{Mapping Route Hotspots:} We map latency and packet loss along the entire route to see if specific spots consistently suffer from poor connectivity, and whether those trouble spots stay the same from day to day.
    
    \item \textbf{Comparing Trips Over Time:} We compare morning versus afternoon runs, inbound versus outbound trips, and week-to-week trends from September through November, and use historical weather records to compare dry and rainy days. This shows how time of day, network traffic, and weather affect performance.
    
    \item \textbf{Tracking Satellite Handovers:} Using logged satellite IDs and GPS, we track every time the terminal switches satellites to see how often handovers trigger latency spikes or temporary dropouts, and whether those handovers are scheduled or forced by blockages.
\end{itemize}

\section{Deliverables}

Below can be found our planned project deliverables and a rough biweekly schedule. Both the deliverables and schedule are subject to change. \\

\textbf{Deliverables:}\\
  - Cleaned/processed dataset (latency + GPS + satellite ID aligned)\\\\
  - Route-based visualizations / performance maps\\\\
  - Handover analysis + latency-correlation results\\\\
  - Statistical/ML model identifying recurring and anomalous patterns\\\\
  - Final report + presentation\\\\
\textbf{Biweekly Schedule:}\\
    - Weeks 1–2: data cleaning and exploration (Completed)\\\\
    - Weeks 3–4: Route-based visualizations / performance maps\\\\
    - Weeks 5–6: Handover analysis + latency-correlation results\\\\
    - Weeks 7–8: Statistical/ML model identifying recurring and anomalous patterns, Final Report\\\\



\bibliographystyle{plain}
\bibliography{references} 


\end{document}